\subsubsection{Heterogeneous Capital Results}
\label{sec:432:specB_results}

Specification A treats aggregate core capital $k_{\text{cap},t} = \log(K_{\text{ME},t} + K_{\text{NRC},t})$ as a single regressor. This collapses the physical structure of the workplace into a scalar, obscuring whether wage pressure operates through the \emph{level} of capital or through its \emph{composition}---the relative weight of machinery and equipment ($K^{\text{ME}}$) versus nonresidential structures ($K^{\text{NRC}}$). Specification B decomposes the aggregate into two components: a structural base ($k^{\text{NRC}}_t = \log K^{\text{NRC}}_t$) and the capital composition ratio ($\kappa_t = k^{\text{ME}}_t - k^{\text{NRC}}_t$), which measures the mechanization gap between equipment and structures in log-levels. The disaggregated cointegrating equation takes the form:
\begin{equation}
	y_t = \alpha + \theta_{\text{NRC}}\, k^{\text{NRC}}_t + \theta_\kappa\, \kappa_t + \theta_{\omega\kappa}\, \xi^{\omega\kappa}_t + \theta_{\kappa p}\, \xi^{p\kappa}_t + \boldsymbol{\gamma}' \mathbf{D}_t + u_t,
	\label{eq:432:specB_model}
\end{equation}
where $\xi^{\omega\kappa}_t$ denotes the FWL-orthogonalized effective wage--mechanization interaction and $\xi^{p\kappa}_t$ denotes the FWL-orthogonalized relative price--mechanization interaction. Both interaction terms are residualized against the two base capital regressors ($k^{\text{NRC}}_t, \kappa_t$) and deterministic step shifts ($\mathbf{D}_t$) prior to estimation. The effective wage interaction under Set~2 is $\xi^{\omega\kappa}_t = \omega_{t-1}(1-\iota_{t-1}) \times \kappa_t$; the relative price interaction is $\xi^{p\kappa}_t = p^{\kappa}_t \times \kappa_t$, where $p^{\kappa}_t$ denotes the relative asset price ratio. Non-standard fixed-$b$ critical values are generated from a dedicated 100,000-pass Monte Carlo simulation engine that replicates the exact Specification~B integration geometry---five independent $I(1)$ stochastic trends, two simultaneous $I(1) \times I(1)$ interaction products, and joint FWL orthogonalization---matching the empirical sample sizes ($T = 94$ and $T = 77$).

\paragraph{Economic content of the decomposition.}
The parameter $\theta_{\text{NRC}}$ captures the long-run capacity yield of structures---factories, warehouses, commercial buildings---which embody extensive accumulation: expanding the physical envelope within which production occurs. The parameter $\theta_\kappa$ captures the additional capacity yield of machinery and equipment \emph{above and beyond} what structures deliver. A positive $\theta_\kappa$ means that a wider mechanization gap---more equipment per dollar of structures---raises productive capacity independently of the total capital stock. The interaction $\theta_{\omega\kappa}$ asks whether effective wage pressure conditions this compositional margin. Under the Okishio-induced mechanization mechanism, rising wages force firms to substitute equipment for labor on the shop floor. If the resulting machinery expansion is defensive---driven by unit-cost containment rather than autonomous technical expansion---it raises the equipment-to-structures ratio without proportionally expanding output capacity, yielding $\theta_{\omega\kappa} < 0$.

\paragraph{Standard wage interactions fail at the intensive margin.}
Tables~\ref{tab:specB_set1_gva_app} and \ref{tab:specB_set1_nva_app} in Appendix~\ref{app:specB_full_results} present the estimation results using the unadjusted wage share across all 24 specifications. The structural base elasticity $\hat{\theta}_{\text{NRC}}$ remains stable and highly significant ($0.47$--$0.61$, $p < 0.01$), confirming that nonresidential structures anchor the cointegrating relationship regardless of wage measure or break treatment. The standalone mechanization gap $\hat{\theta}_\kappa$ is positive in nearly all models ($0.34$--$1.11$), indicating that equipment contributes to capacity above structures. The wage--mechanization interaction $\hat{\theta}_{\omega\kappa}$, however, fails to achieve statistical significance in any $T = 94$ model, with point estimates ranging from $-4.89$ ($t = -1.58$) to $+7.95$ ($t = 1.33$). The standard wage interaction fails at the intensive margin for the same reason it fails at the aggregate level in Specification~A: unadjusted wage shares conflate supervisory overhead with productive variable capital and ignore the post-Fordist diversion of investment toward intellectual property products.

\paragraph{Effective wage pressure and defensive mechanization at the intensive margin.}
Table~\ref{tab:specB_set2_nva_main} presents the main estimates using the effective wage pressure measure for the NVA basis. When effective wage pressure is scaled by $(1-\iota_{t-1})$ and interacted with the mechanization gap $\kappa_t$, the wage--mechanization interaction $\hat{\theta}_{\omega\kappa}$ is negative and statistically significant across all eight models. In the preferred Clean Corporate Tier~1 NVA specification with Bai-Perron step shifts (Column~4), $\hat{\theta}_{\omega\kappa} = -7.93$ ($SE = 2.96$, $p < 0.05$), with a structural base elasticity of $\hat{\theta}_{\text{NRC}} = 0.60$ ($SE = 0.08$, $p < 0.01$) and a standalone mechanization gap of $\hat{\theta}_\kappa = 0.39$ ($SE = 0.29$). In the Shaikh-Tonak Tier~2 NVA specification without breaks (Column~7), the interaction strengthens to $\hat{\theta}_{\omega\kappa} = -17.05$ ($SE = 2.34$, $p < 0.01$), with $\hat{\theta}_{\text{NRC}} = 0.54$ ($SE = 0.01$) and $\hat{\theta}_\kappa = 0.33$ ($SE = 0.06$, $p < 0.01$). The Clean Corporate Tier~2 NVA models (Columns~5--6) yield $\hat{\theta}_{\omega\kappa} = -14.97$ ($SE = 2.01$, $p < 0.01$), showing the tightest standard errors in the table.

The negative sign of $\hat{\theta}_{\omega\kappa}$ provides direct evidence of defensive mechanization operating at the \emph{intensive margin} of capital composition. Specification~A established that effective wage pressure depresses the aggregate capacity yield of total capital ($\hat{\theta}_1 = -0.30$). Specification~B isolates the channel: wage pressure forces firms to widen the mechanization gap---substituting equipment for structures on the shop floor---but the capacity yield of that compositional shift is negative. The firm accumulates more machinery per dollar of structures, yet productive capacity does not expand proportionally, because the substitution is driven by cost containment rather than by genuine technical complementarity between new equipment and existing plant.\footnote{The apparent difference in magnitude between Specification~A ($\hat{\theta}_1 \approx -0.30$) and Specification~B ($\hat{\theta}_{\omega\kappa} \approx -7.93$ to $-17.05$) reflects regressor scaling, not economic inconsistency. In Specification~A, $\theta_1$ multiplies $\omega_{t-1}(1-\iota_{t-1}) \times k_{\text{cap},t}$, where aggregate capital $k_{\text{cap},t} \approx 16$--$17$ in log-levels. In Specification~B, $\theta_{\omega\kappa}$ multiplies $\omega_{t-1}(1-\iota_{t-1}) \times \kappa_t$, where the mechanization gap $\kappa_t = k^{\text{ME}}_t - k^{\text{NRC}}_t \approx 0.3$--$0.8$. The total marginal effect (coefficient $\times$ regressor scale) is comparable across both specifications.}

\paragraph{Relative price interactions.}
The relative asset price interaction $\hat{\theta}_{\kappa p}$ is generally small and statistically insignificant. In the NFC Baseline and Clean Corporate Tier~1 models (Columns~1--4), it is statistically indistinguishable from zero. In the Shaikh-Tonak Tier~2 specification (Columns~7--8), $\hat{\theta}_{\kappa p} = -0.79$ ($SE = 0.32$, $p < 0.05$), suggesting that when equipment becomes relatively more expensive, the capacity yield of the mechanization gap declines. However, this result is fragile and does not hold across wage share measures. The Okishio mechanism operates through distributional conflict over newly created value, not through relative asset prices.

\paragraph{Cointegration verdict.}
The effective wage models show strong evidence of cointegration. The Shin LM test never rejects the null of cointegration ($C_{\text{Shin}} \le 0.03$, $p > 0.10$), and the Vogelsang-Wagner Type~2 Wald statistic remains well below critical thresholds ($S_{\text{Type~2}} \le 15.37$). The Engle-Granger residual ADF test rejects the null of no cointegration at the 5\% level in every break-adjusted model: $t_{\text{ADF}} = -4.58$ (Column~2), $-4.60$ (Column~4), $-4.71$ (Column~7), and $-4.57$ (Column~8). Without step shifts, the ADF test rejects at the 10\% level across all $T = 94$ models ($t_{\text{ADF}} \le -4.33$) and at the 5\% level for the $T = 76$ Shaikh-Tonak Tier~2 model ($t_{\text{ADF}} = -4.71$, Column~7). The full cointegration triad---Shin, Vogelsang-Wagner, and ADF---confirms that the disaggregated capital composition, real output, and effective wage pressure form a stable long-run relationship.

\paragraph{Robustness.}
Including Bai-Perron step shifts absorbs historical level discontinuities, but leaves the sign and significance of the wage--mechanization interaction unchanged. In the Clean Corporate Tier~1 NVA model, controlling for breaks at 1944, 1977, and 2001 sharpens the ADF rejection from $-4.34$ ($p < 0.10$) to $-4.60$ ($p < 0.05$) while reducing the Vogelsang-Wagner statistic from $15.37$ to $6.37$. The wage interaction $\hat{\theta}_{\omega\kappa}$ moves from $-11.08$ (Column~3) to $-7.93$ (Column~4). This modest attenuation indicates that the step dummies absorb part of the interaction's raw variance. The interaction remains negative and significant at the 5\% level, confirming that the compositional conditioning channel survives institutional break controls.

The structural base elasticity $\hat{\theta}_{\text{NRC}}$ falls consistently within the $0.51$--$0.61$ band across wage share measures, closely tracking the aggregate $\hat{\theta}_0 = 0.58$ from Specification~A. This implies the capacity yield of structures is a stable parameter of U.S.\ capital accumulation. The wage interaction $\hat{\theta}_{\omega\kappa}$ grows larger in absolute value as we remove supervisory overhead from the wage share: from $-10.13$ in the NFC Baseline (Column~2) to $-7.93$ in Clean Corporate Tier~1 (Column~4) to $-16.76$ in Shaikh-Tonak Tier~2 (Column~8). Removing managerial salaries isolates the productive variable capital driving the firm's cost-minimization decision, which sharpens the defensive mechanization effect.

\input{tables/tab_specB_set2_nva_main.tex}

\paragraph{The aggregate transformation elasticity along the historical path.}
Having established the structural asset-weighted transformation elasticity $\theta^{\text{Aggregate}}_t$ in Equation~\eqref{eq:32:theta_aggregate_structural} (Section~3.2), I evaluate its empirical trajectory using the estimated parameters from the preferred Specification~B model (Clean Corporate Tier~1 NVA, Set~2, Column~4 of Table~\ref{tab:specB_set2_nva_main}). The structural marginal capacity yield of mechanization $MP_{\kappa,t}^{\text{str}} = \beta_\kappa + \theta_{\omega\kappa} C_{\text{index},t} + \theta_{\kappa p} p^\kappa_t$ is driven directly by the state-dependent levels of effective wage pressure ($C_{\text{index},t}$) and relative asset prices ($p^\kappa_t$), where initial-observation centering ($\theta_\kappa^{\text{str}} = 0.39$) serves as an anchoring convention to identify the baseline dimensionless yield at $t=1931$.

Figure~\ref{fig:specB_candidate1} plots $\theta^{\text{Aggregate}}_t$ against the dynamic Specification~A baseline ($\theta^{\text{SpecA}}_t = 0.58 - 0.30 \cdot C_{\text{index},t}$) and the Harrodian knife-edge threshold ($\theta = 1.00$) under standard structural regime shading.

\begin{figure}[h!]
	\centering
	\caption{Aggregate Transformation Elasticity ($\theta^{\text{Aggregate}}_t$), Specification~B vs.\ Baseline}
	\label{fig:specB_candidate1}
	\includegraphics[width=0.88\textwidth]{fig_specB_candidate1_aggregate_elasticity.png}
	\begin{tabular}{p{0.88\textwidth}}
		\hline
		\footnotesize \textbf{Notes:} Aggregate transformation elasticity $\theta^{\text{Aggregate}}_t$ (Navy line) computed according to Equation~\eqref{eq:32:theta_aggregate_structural} using FWL structural parameters from the preferred Specification~B model (Clean Corporate Tier~1 NVA, Set~2, Column~4 of Table~\ref{tab:specB_set2_nva_main}). The Slate Blue line traces the dynamic Specification~A baseline ($\theta^{\text{SpecA}}_t = 0.58 - 0.30 \, C_{\text{index},t}$). The dotted horizontal line at $\theta = 0.58$ marks the constant base from Specification~A. The dashed line at $\theta = 1.00$ marks the Harrodian knife-edge. Background shading indicates structural accumulation regimes. \\
		\hline
	\end{tabular}
\end{figure}

\paragraph{Empirical comparison: Specification~B disciplines the Specification~A baseline.}
Figure~\ref{fig:specB_candidate1} delivers two core empirical findings. First, the structural aggregate elasticity $\theta^{\text{Aggregate}}_t$ closely tracks the dynamic Specification~A path but operates within a narrower, disciplined band---$[0.32, 0.86]$ versus $[0.37, 0.75]$. The two paths move together because both capture the same underlying macroeconomic force: effective wage pressure depresses the marginal capacity yield of capital across both specifications, and the estimated base elasticity of structures ($\hat{\theta}_{\text{NRC}} \approx 0.60$) matches Specification~A's aggregate base ($\hat{\theta}_0 = 0.58$). Second, the Specification~B trajectory remains strictly below the Harrodian knife-edge ($\theta < 1$) across the entire 1931--2024 sample, confirming that aggregate capacity formation in the U.S.\ economy operates permanently within the sub-Harrodian stability regime.

The disaggregated capital structure explains what Specification~A's aggregate elasticity averages over. During the Fordist era ($C_{\text{index},t} > 0.65$), high effective wage pressure drove negative marginal mechanization yields ($MP_{\kappa,t}^{\text{str}} < 0$). Firms substituted equipment for labor to contain unit labor costs, but because this expansion was defensive rather than technically autonomous, it depressed the aggregate capacity yield of total capital. In the Neoliberal period ($C_{\text{index},t} < 0.50$), declining wage pressure restored $MP_{\kappa,t}^{\text{str}}$ to positive territory, pulling $\theta^{\text{Aggregate}}_t$ back toward the structural base. Specification~A collapses this compositional mechanism into a scalar interaction ($\hat{\theta}_1 = -0.30$); Specification~B opens the intensive margin and confirms the Okishio mechanism empirically.

\begin{figure}[h!]
	\centering
	\caption{Empirical Accumulation Path Derivatives ($dk^{\text{ME}} / dk^{\text{NRC}}$ and $d\kappa / dk^{\text{NRC}}$)}
	\label{fig:specB_dkappa_dknrc_slope}
	\includegraphics[width=0.88\textwidth]{fig_specB_dkappa_dknrc_slope.png}
	\begin{tabular}{p{0.88\textwidth}}
		\hline
		\footnotesize \textbf{Notes:} Empirical accumulation path derivatives estimated via non-parametric LOESS time-domain chain rule ($\text{span} = 0.35$). The green solid line measures the relative accumulation slope $dk^{\text{ME}}_t / dk^{\text{NRC}}_t$. The red dot-dashed line measures the mechanization gap derivative $d\kappa_t / dk^{\text{NRC}}_t = (dk^{\text{ME}}_t / dk^{\text{NRC}}_t) - 1$. The horizontal dashed line at $1.00$ indicates balanced expansion where equipment and structures grow at identical rates. \\
		\hline
	\end{tabular}
\end{figure}

\paragraph{Empirical failure of the conventional Divisia derivative path.}
Figure~\ref{fig:specB_dkappa_dknrc_slope} illustrates why the empirical total-derivative path $\theta^{\text{Aggregate}}_{\text{total},t}$ (Equation~\eqref{eq:32:theta_aggregate}) fails in practice. In the mid-1990s, corporate accumulation shifted away from physical machinery toward intellectual property products and financial assets, causing physical equipment growth to slow relative to structures ($dk^{\text{ME}}_t / dk^{\text{NRC}}_t < 1.00$). As $d\kappa_t / dk^{\text{NRC}}_t$ crossed zero, the Divisia denominator in Equation~\eqref{eq:32:theta_aggregate} approached zero, sending the raw total-derivative elasticity to spurious infinite spikes. As established in Section~3.2, this instability reflects the theoretical pathology of differentiating a value aggregate under structural accumulation shifts. The structural asset-weighted elasticity $\theta^{\text{Aggregate}}_t$ in Equation~\eqref{eq:32:theta_aggregate_structural} purges these Divisia singularities, yielding a continuous, empirically well-behaved trajectory.

